\chapter{Gaussian Elimination}

In the previous chapter, we saw linear systems $\mat{A}\vec{x}=\vec{b}$ through
two geometric lenses: intersecting hyperplanes in the row picture, and combining
column arrows in the column picture. Both tell us what a solution means, but neither
gives a practical way to compute it when a system has ten or a thousand equations.

We need a systematic method: \emph{Gaussian elimination}. The strategy is
simple: add multiples of one equation to another to eliminate variables until
the system is triangular, then solve from the bottom up by back-substitution.
Start, as always, with a picture.

\section{The picture: collapsing planes onto a corner}

Recall the two geometric faces of a linear system. In the row picture, each
equation is a constraint that cuts down the space of candidates. In two dimensions,
two non-parallel lines cross at a single point; in three dimensions, three tilted
planes meet at a point. Solving the system means \emph{pinning down that shared point}.

Gaussian elimination finds that point through a direct strategy:
\textbf{we replace the tilted equations with simpler, axis-aligned equations
that have the exact same solution, until the coordinates can be read off one by one.}

In the row picture, each step gently swings a tilted plane until it becomes parallel
to a coordinate plane, without ever budging the intersection point. In the column
picture, that same step introduces zeros into the column vectors, uncoupling the
variables so that the final column alone accounts for the bottom coordinate of
$\vec{b}$, letting us solve the weights in reverse order.

\begin{intuitionbox}
Imagine three tilted planes in space crossing at a single hidden point. You
cannot see where they meet because the planes are at awkward angles. Gaussian
elimination gently rotates and tilts the equations -- without ever moving the
crossing point -- until one plane becomes perfectly horizontal (it fixes $z$ outright),
the next pins down $y$ once $z$ is known, and the last pins down $x$ once $y$ and
$z$ are known. The crossing point never budges. We have only changed our
\emph{description} of it, from three vague clues into three sharp ones.
\end{intuitionbox}

The key guarantee is that the crossing point never moves. Everything we do
will be a \emph{solution-preserving} operation: it changes the equations but not
their common solution. That claim is the foundation everything else rests on, so
let us pin it down.

Figure~\ref{ch07:two-lines} shows the simplest case -- two lines in the plane --
before and after one elimination step. The replacement move swings the second
line into a horizontal one (it now constrains only $y$), and the crossing point
stays exactly where it was.

\begin{figure}[H]
\centering
\begin{tikzpicture}[scale=0.9]
\begin{scope}
  \draw[->,gray] (-0.4,-0.4)--(3.4,-0.4) node[right]{\small$x$};
  \draw[->,gray] (-0.4,-0.4)--(-0.4,3.0) node[above]{\small$y$};
  % two tilted lines crossing near (1.5,1.2)
  \draw[thick,vecblue] (-0.2,2.4)--(3.2,0.05);
  \draw[thick,vecred] (0.0,-0.3)--(2.7,2.4);
  \fill[black] (1.5,1.2) circle (2pt);
  \node[right=4pt] at (1.5,1.2){\small solution $(x^*,y^*)$};
  \node[below] at (1.4,-0.9){\small two tilted lines};
\end{scope}
\begin{scope}[xshift=6cm]
  \draw[->,gray] (-0.4,-0.4)--(3.4,-0.4) node[right]{\small$x$};
  \draw[->,gray] (-0.4,-0.4)--(-0.4,3.0) node[above]{\small$y$};
  \draw[thick,vecblue] (-0.2,2.4)--(3.2,0.05);
  % second line now horizontal through y=1.2 (pins y)
  \draw[thick,vecred] (-0.3,1.2)--(3.2,1.2);
  \fill[black] (1.5,1.2) circle (2pt);
  \node[below left=3pt] at (1.5,1.2){\small same point $(x^*,y^*)$};
  \draw[->,vecred!70] (2.7,2.3) to[bend right=20] (2.8,1.35);
  \node[below] at (1.4,-0.9){\small after elimination: red pins $y$};
\end{scope}
\end{tikzpicture}
\caption{Elimination in the plane. One replacement step swings the red line into
a horizontal one that fixes $y$ outright, while the crossing point (the solution)
does not move. Back-substitution then reads $x$ off the blue line.}
\label{ch07:two-lines}
\end{figure}

\subsection{The three legal moves}

We are allowed exactly three operations on the equations of a system. Each one
is reversible, and each one leaves the set of solutions completely unchanged.

\begin{focusbox}
\textbf{The three elementary row operations.}
\begin{enumerate}
  \item \textbf{Scaling:} multiply an equation by a nonzero number.
  \item \textbf{Replacement:} add a multiple of one equation to another equation.
  \item \textbf{Exchange:} swap the positions of two equations.
\end{enumerate}
Every one of these is reversible, so none of them can lose or invent solutions.
\end{focusbox}

Why do these preserve the solution set? Consider row replacement. Suppose
a point $(x,y,z)$ satisfies equation $E_1$ and equation $E_2$. Then it certainly
satisfies $E_2 + cE_1$ for any number $c$, because both sides of $E_1$ are equal
at that point, so adding $c$ times a true equation to another true equation
keeps it true. Conversely, if a point satisfies $E_1$ and $E_2 + cE_1$, we can
\emph{recover} $E_2$ as $(E_2 + cE_1) - cE_1$, so the point satisfies $E_2$ too.
The two systems have identical solution sets. Scaling by a nonzero number and
swapping are reversible for the obvious reasons. (Scaling by zero is forbidden
precisely because it is \emph{not} reversible -- it would erase an equation.)

\begin{trapbox}
The forbidden move is multiplying an equation by \textbf{zero}. ``$0 = 0$'' is
true everywhere, so it throws away the information that equation carried, and you
can never get it back. Every legal move must be undoable. Keep that test in your
head: if you could not reverse it, you may not do it.
\end{trapbox}

\subsection{From equations to a grid of numbers}

Writing $x$, $y$, $z$ and plus signs over and over is a waste of ink. The
variables just hold places; what matters are the coefficients. So we strip the
system down to its numbers and stack them in a rectangular grid. The system
\[
\begin{aligned}
2x + 1y + 1z &= 5\\
4x + 1y + 0z &= 6\\
-2x + 2y + 1z &= 1
\end{aligned}
\qquad\longleftrightarrow\qquad
\left[\begin{array}{ccc|c}
2 & 1 & 1 & 5\\
4 & 1 & 0 & 6\\
-2 & 2 & 1 & 1
\end{array}\right].
\]
The vertical bar separates the coefficient matrix $\mat{A}$ from the right-hand
side $\vec{b}$. The whole grid is the \emph{augmented matrix}, written
$[\,\mat{A}\mid\vec{b}\,]$. Each row is one equation; each of the first three
columns belongs to one variable. The three row operations now become operations
on rows of numbers -- which is exactly what a computer wants.

\section{Elimination to upper-triangular form}

The goal of the elimination phase is to make the coefficient matrix
\emph{upper triangular}: every entry below the main diagonal becomes zero. Why
that shape? Because an upper-triangular system is a staircase you can climb from
the bottom. The last equation involves only the last variable, so you solve it
immediately. The second-to-last involves only the last two, and you already know
one of them. And so on up to the top. This climb is called
\emph{back-substitution}, and it solves line by line without guesswork once the triangle is built.
Figure~\ref{ch07:triangle} shows the shape we are driving toward: zeros below the
diagonal, pivots on it.

\begin{figure}[H]
\centering
\begin{tikzpicture}[scale=0.62]
  % full matrix
  \draw[thick,sourceblue] (0,0) rectangle (3,3);
  \foreach \i in {1,2} \draw[gray!40] (0,\i)--(3,\i);
  \foreach \i in {1,2} \draw[gray!40] (\i,0)--(\i,3);
  \node at (0.5,2.5){$*$}; \node at (1.5,2.5){$*$}; \node at (2.5,2.5){$*$};
  \node at (0.5,1.5){$*$}; \node at (1.5,1.5){$*$}; \node at (2.5,1.5){$*$};
  \node at (0.5,0.5){$*$}; \node at (1.5,0.5){$*$}; \node at (2.5,0.5){$*$};
  \node[below] at (1.5,-0.3){full matrix $\mat{A}$};
  % arrow
  \draw[-{Stealth[length=3mm]},thick] (3.6,1.5)--(5.1,1.5);
  \node[above] at (4.35,1.6){\small eliminate};
  % triangular
  \begin{scope}[xshift=5.7cm]
  \draw[thick,targetgreen] (0,0) rectangle (3,3);
  \foreach \i in {1,2} \draw[gray!40] (0,\i)--(3,\i);
  \foreach \i in {1,2} \draw[gray!40] (\i,0)--(\i,3);
  \fill[red!10] (0,0) -- (2,0) -- (2,1) -- (1,1) -- (1,2) -- (0,2) -- cycle;
  \node[structurepurple] at (0.5,2.5){$p$}; \node at (1.5,2.5){$*$}; \node at (2.5,2.5){$*$};
  \node[trapred] at (0.5,1.5){$0$}; \node[structurepurple] at (1.5,1.5){$p$}; \node at (2.5,1.5){$*$};
  \node[trapred] at (0.5,0.5){$0$}; \node[trapred] at (1.5,0.5){$0$}; \node[structurepurple] at (2.5,0.5){$p$};
  \node[below] at (1.5,-0.3){upper triangular $\mat{U}$};
  \end{scope}
\end{tikzpicture}
\caption{Elimination drives every entry below the diagonal to zero (shaded
region). The surviving diagonal entries marked $p$ are the \emph{pivots}.}
\label{ch07:triangle}
\end{figure}

\subsection{Pivots and multipliers}

The mechanics are simple. We work column by column, left to right. In each
column we pick a \emph{pivot}: the diagonal entry we will use to knock out
everything below it. To clear an entry below the pivot, we subtract the right
multiple of the pivot's row.

\begin{focusbox}
To eliminate the entry in row $i$, column $j$, using the pivot in row $j$,
column $j$, the \textbf{multiplier} is
\[
\ell_{ij} \;=\; \frac{\text{entry to kill}}{\text{pivot}}
\;=\; \frac{a_{ij}}{a_{jj}},
\]
and the operation is: \emph{new row $i$} $=$ \emph{old row $i$}
$-\;\ell_{ij}\times(\text{pivot row } j)$. Choosing $\ell_{ij}$ this way makes
the target entry become exactly $a_{ij} - \ell_{ij}\,a_{jj} = a_{ij} - a_{ij} =
0$. Hold onto these multipliers; they encode the entire elimination and will
reappear in Chapter~\ref{ch09:triangular-picture} as the entries of the matrix $\mat{L}$.
\end{focusbox}

\begin{recoverybox}
\textbf{Pause and think:} Suppose you reach the second row and find a zero in its pivot position. Could you fix it by adding a multiple of row 1? What would happen to column 1 if you did?
\end{recoverybox}

The answer is no: adding row 1 would dump nonzero entries right back into column 1 of row 2, destroying the zero just created. Previously computed zeros must be preserved. That is why the algorithm always looks \emph{below} for pivots, never above.

\subsection{A 3x3 traced fully}

Let us run the example above all the way through. Watch every multiplier.
\[
\left[\begin{array}{ccc|c}
2 & 1 & 1 & 5\\
4 & 1 & 0 & 6\\
-2 & 2 & 1 & 1
\end{array}\right].
\]
The first pivot is the $2$ in the top-left corner. We use it to clear the rest of
column 1.

\textbf{Step 1.} Kill the $4$ in row 2. Multiplier $\ell_{21} = 4/2 = 2$.
Replace row 2 with (row 2) $- 2\cdot$(row 1):
\[
[\,4,\;1,\;0\mid 6\,] - 2[\,2,\;1,\;1\mid 5\,]
= [\,0,\;-1,\;-2\mid -4\,].
\]
\[
\left[\begin{array}{ccc|c}
2 & 1 & 1 & 5\\
0 & -1 & -2 & -4\\
-2 & 2 & 1 & 1
\end{array}\right].
\]

\textbf{Step 2.} Kill the $-2$ in row 3. Multiplier $\ell_{31} = -2/2 = -1$.
Replace row 3 with (row 3) $-(-1)\cdot$(row 1) $=$ (row 3) $+$ (row 1):
\[
[\,-2,\;2,\;1\mid 1\,] + [\,2,\;1,\;1\mid 5\,] = [\,0,\;3,\;2\mid 6\,].
\]
\[
\left[\begin{array}{ccc|c}
2 & 1 & 1 & 5\\
0 & -1 & -2 & -4\\
0 & 3 & 2 & 6
\end{array}\right].
\]
Column 1 is clean. Now we move to the second pivot, the $-1$ now sitting in row
2, column 2, and use it to clear column 2 below.

\textbf{Step 3.} Kill the $3$ in row 3. Multiplier $\ell_{32} = 3/(-1) = -3$.
Replace row 3 with (row 3) $-(-3)\cdot$(row 2) $=$ (row 3) $+ 3\cdot$(row 2):
\[
[\,0,\;3,\;2\mid 6\,] + 3[\,0,\;-1,\;-2\mid -4\,] = [\,0,\;0,\;-4\mid -6\,].
\]
\[
\mat{U}=\left[\begin{array}{ccc|c}
2 & 1 & 1 & 5\\
0 & -1 & -2 & -4\\
0 & 0 & -4 & -6
\end{array}\right].
\]
We have reached upper-triangular form. The three pivots are $2,\,-1,\,-4$, all
sitting on the diagonal. The forward sweep is done.

\subsection{Back-substitution: climbing the staircase}

Now read the triangle from the bottom up. The last row says
\[
-4z = -6 \quad\Longrightarrow\quad z = \tfrac{3}{2}.
\]
The middle row says $-y - 2z = -4$. Plug in $z = 3/2$:
\[
-y - 2\cdot\tfrac{3}{2} = -4 \;\Longrightarrow\; -y - 3 = -4 \;\Longrightarrow\;
-y = -1 \;\Longrightarrow\; y = 1.
\]
The top row says $2x + y + z = 5$. Plug in $y = 1$ and $z = 3/2$:
\[
2x + 1 + \tfrac{3}{2} = 5 \;\Longrightarrow\; 2x = \tfrac{5}{2}
\;\Longrightarrow\; x = \tfrac{5}{4}.
\]
So the three tilted planes meet at the single point
$\bigl(\tfrac54,\,1,\,\tfrac32\bigr)$. You should always sanity-check the
original equations. The first: $2\cdot\frac54 + 1 + \frac32 = \frac52 + 1 +
\frac32 = 5$. Good. The second: $4\cdot\frac54 + 1 + 0 = 5 + 1 = 6$. Good. The
third: $-2\cdot\frac54 + 2\cdot 1 + \frac32 = -\frac52 + 2 + \frac32 = 1$. Good.

Figure~\ref{ch07:back-sub} shows this triangular climb in action. Once the zeros
are locked in place below the diagonal, the bottom equation hands you the last
variable immediately. Substituting that value upward determines the next variable, and so
on up the triangular system until the entire solution vector is assembled.

\begin{figure}[H]
\centering
\begin{tikzpicture}[scale=0.95, >=Stealth]
  \node[draw=sourceblue, rounded corners, fill=blue!4!white, inner sep=6pt, text width=7.4cm] (row3) at (0, 0) {
    \textbf{Step 1 (Row 3):} $-4z = -6 \implies \mathbf{z = \tfrac{3}{2}}$
  };
  \node[draw=sourceblue, rounded corners, fill=blue!4!white, inner sep=6pt, text width=7.4cm] (row2) at (0, 1.8) {
    \textbf{Step 2 (Row 2):} $-y - 2(\mathbf{\tfrac{3}{2}}) = -4 \implies \mathbf{y = 1}$
  };
  \node[draw=sourceblue, rounded corners, fill=blue!4!white, inner sep=6pt, text width=7.4cm] (row1) at (0, 3.6) {
    \textbf{Step 3 (Row 1):} $2x + \mathbf{1} + \mathbf{\tfrac{3}{2}} = 5 \implies \mathbf{x = \tfrac{5}{4}}$
  };

  \draw[->, thick, targetgreen] (-4.3, -0.2) -- (-4.3, 3.8) node[midway, left=2pt, targetgreen, align=center] {\small Climb upward\\\small (back-substitute)};

  \draw[->, thick, quotientorange] (row3.north) to[bend right=25] node[right=2pt, text=quotientorange]{\small substitute $z$} (row2.south);
  \draw[->, thick, quotientorange] (row2.north) to[bend right=25] node[right=2pt, text=quotientorange]{\small substitute $y, z$} (row1.south);

  \node[draw=structurepurple, rounded corners, fill=purple!5!white, inner sep=5pt, right=0.8cm of row2] (sol) {
    $\vec{x} = \colvec{\tfrac{5}{4} \\[3pt] 1 \\[3pt] \tfrac{3}{2}}$
  };
  \draw[->, thick, structurepurple] (row1.east) to[bend left=20] (sol.north);
\end{tikzpicture}
\caption{Back-substitution climbs the triangular system from the bottom up.
Row 3 fixes $z$ immediately. Feeding $z=\frac{3}{2}$ into Row 2 isolates $y=1$.
Feeding both into Row 1 isolates $x=\frac{5}{4}$, assembling the complete solution vector $\vec{x}$.}
\label{ch07:back-sub}
\end{figure}

\begin{focusbox}
\textbf{The two phases.} \emph{Forward elimination} reduces the matrix to an
upper triangle $\mat{U}$, recording multipliers as it goes. \emph{Back-
substitution} then climbs the triangle from the bottom, solving one variable at
each step. Together they solve the system. Forward elimination is the expensive
part; back-substitution requires fewer operations.
\end{focusbox}

\begin{strangbox}
Strang opens his course with exactly this: ``the geometry of linear equations''
followed immediately by ``elimination with matrices.'' If you are reading this
\emph{before} Strang, you now have the geometric view -- planes pinned down one
coordinate at a time -- that clarifies the reasoning behind each algebraic step.
If you are reading \emph{after} Strang, notice that we are being deliberate about
the multipliers $\ell_{ij}$; they are recorded systematically rather than discarded.
They are the basis of the $\mat{A}=\mat{LU}$ factorization that Chapter~\ref{ch09:triangular-picture}
makes precise, and they encode the entire elimination as a single matrix.
\end{strangbox}

\begin{labbox}
Solve this $2\times 2$ system by hand, tracing both phases:
\[
\begin{aligned}
3x + 2y &= 7\\
6x + 5y &= 16.
\end{aligned}
\]
Pivot on the $3$, find the multiplier $6/3 = 2$, eliminate, then back-substitute.
You should get $y = 2$, $x = 1$. Check it against both original equations before
moving on -- the checking habit is half the skill.
\end{labbox}

\section{Counting the work, and why triangles are cheap}

Why is this the algorithm everyone actually uses? Cost. For an $n\times n$ system, forward
elimination does on the order of $\tfrac13 n^3$ arithmetic operations -- one
nested loop over rows, columns, and the entries to update. Back-substitution
does only about $\tfrac12 n^2$. The forward sweep dominates completely, as
Figure~\ref{ch07:cost} illustrates.

\begin{figure}[H]
\centering
\begin{tikzpicture}[scale=0.95]
\begin{axis}[
  width=10.5cm, height=6.6cm,
  xlabel={system size $n$}, ylabel={arithmetic operations},
  domain=1:50, samples=80,
  legend pos=north west, legend cell align=left,
  grid=both, grid style={gray!18},
  xmin=0, xmax=50, ymin=0,
  axis lines=left, tick align=outside,
  every axis plot/.append style={very thick},
]
  \addplot[sourceblue]{(1/3)*x^3};
  \addlegendentry{forward elimination $\sim\tfrac13 n^3$}
  \addplot[targetgreen]{(1/2)*x^2};
  \addlegendentry{back-substitution $\sim\tfrac12 n^2$}
\end{axis}
\end{tikzpicture}
\caption{The two phases compared. The cubic forward sweep (blue) outruns the
quadratic back-substitution (green) so fast that for any sizeable system the
climb down the triangle is essentially free. Doubling $n$ multiplies the blue
curve by eight.}
\label{ch07:cost}
\end{figure}

\begin{aestheticbox}
That $\tfrac13 n^3$ is one of the most consequential numbers in computing.
Doubling the size of your system makes elimination eight times slower, not twice
-- which is why large sparse systems (where most coefficients are zero) require
special algorithms that avoid filling in those zeros. But for dense systems,
standard elimination remains, asymptotically, among the most effective algorithms
for general use.
\end{aestheticbox}

\begin{connectionbox}
The $n^3$ scaling is the same wall you hit when multiplying two matrices the
naive way, and it is no coincidence. Solving $\mat{A}\vec{x}=\vec{b}$,
inverting $\mat{A}$, and multiplying matrices are all $O(n^3)$ problems that turn
out to be equivalent in difficulty up to constants -- a fact from the theory of
algorithmic complexity. Calculus students will recognize the family resemblance
to how the cost of a triple integral grows with resolution: each new dimension
multiplies the work. Elimination is what high-dimensional bookkeeping costs.
\end{connectionbox}

\section{Row exchanges: when the pivot is in the wrong place}

So far every pivot we needed was conveniently nonzero and sitting on the
diagonal. Reality is not always so kind. Sometimes a zero shows up exactly where
we wanted a pivot. Consider
\[
\left[\begin{array}{ccc|c}
0 & 2 & 1 & 4\\
1 & 1 & 1 & 6\\
2 & 1 & 3 & 11
\end{array}\right].
\]
The top-left entry is $0$. We cannot use it as a pivot -- the multiplier would
be undefined (division by zero). But this is easily resolved. The equations are
simply listed in an order that placed a zero on the diagonal. We exchange
two rows to bring a nonzero number into the pivot position. Exchange
row 1 and row 2:
\[
\left[\begin{array}{ccc|c}
1 & 1 & 1 & 6\\
0 & 2 & 1 & 4\\
2 & 1 & 3 & 11
\end{array}\right].
\]
Now the pivot is $1$, and elimination proceeds. Clear column 1: row 3 already
has a $2$ to remove, multiplier $\ell_{31}=2/1=2$, so row 3 becomes (row 3)
$-2\cdot$(row 1) $= [\,0,\,-1,\,1\mid -1\,]$. Row 2 already starts with $0$, so
nothing to do there.
\[
\left[\begin{array}{ccc|c}
1 & 1 & 1 & 6\\
0 & 2 & 1 & 4\\
0 & -1 & 1 & -1
\end{array}\right].
\]
Second pivot is $2$ in row 2. Clear below: multiplier $\ell_{32}=-1/2$, so row 3
becomes (row 3) $-(-\tfrac12)\cdot$(row 2) $=$ (row 3) $+\tfrac12\cdot$(row 2)
$= [\,0,\,0,\,\tfrac32\mid 1\,]$.
\[
\left[\begin{array}{ccc|c}
1 & 1 & 1 & 6\\
0 & 2 & 1 & 4\\
0 & 0 & \tfrac32 & 1
\end{array}\right].
\]
Back-substitution: $\tfrac32 z = 1$ gives $z = \tfrac23$; then $2y + \tfrac23 =
4$ gives $y = \tfrac53$; then $x + \tfrac53 + \tfrac23 = 6$ gives $x = \tfrac{11}{3}$.

\begin{focusbox}
A zero in the pivot position is \textbf{not} necessarily a sign of trouble. If
there is a nonzero entry \emph{below} it in the same column, just swap rows to
bring that entry up, and carry on. The swap costs nothing and changes no
solution -- you are only reordering the equations.
\end{focusbox}

\begin{applicationbox}
Computers exchange rows even when they do not strictly have to, and the reason is
\emph{numerical stability} -- keeping tiny, inevitable rounding errors in computer
arithmetic from blowing up into completely wrong answers. Dividing by a pivot that
is tiny -- say $0.0001$ -- multiplies any rounding error in that row by ten thousand,
and that magnified noise can quickly swamp the true answer.

So practical solvers use \emph{partial pivoting}: before clearing a column, scan all
the available rows on or below the diagonal, and swap the row with the largest
absolute entry into the pivot position. This guarantees that every multiplier satisfies
$|\ell_{ij}| \le 1$, ensuring that row subtractions never amplify errors. Every time
\texttt{numpy.linalg.solve}, MATLAB's backslash, or a LAPACK routine solves a system,
partial pivoting runs automatically. It connects textbook math
with exact fractions and scientific software working in finite floating-point arithmetic.
\end{applicationbox}

\section{What goes wrong: zero pivots that will not move}

Row exchange rescues us when a pivot position holds a zero but the column below
has something nonzero to swap in. What if the \emph{entire} column below (and
including) the pivot position is zero? Then no swap can produce a pivot. This is
not merely a bookkeeping detail; it indicates a structural property of the system.
The geometry has changed.

\subsection{Case one: no solution (an inconsistent system)}

Consider
\[
\begin{aligned}
x + 2y &= 3\\
2x + 4y &= 7.
\end{aligned}
\qquad
\left[\begin{array}{cc|c}
1 & 2 & 3\\
2 & 4 & 7
\end{array}\right].
\]
Pivot on the $1$. Multiplier $\ell_{21}=2$. Row 2 becomes (row 2) $-2\cdot$(row
1) $= [\,0,\,0\mid 1\,]$.
\[
\left[\begin{array}{cc|c}
1 & 2 & 3\\
0 & 0 & 1
\end{array}\right].
\]
The second row now reads $0x + 0y = 1$, that is, $0 = 1$. There is no point in
the plane where $0$ equals $1$. The system has \textbf{no solution}.
Geometrically the two lines $x+2y=3$ and $2x+4y=7$ are \emph{parallel}: same
slope, different intercepts, never crossing. Elimination found this for us
automatically -- a row of zeros on the left with a nonzero on the right is the
algebraic fingerprint of parallel, non-intersecting geometry. The left panel of
Figure~\ref{ch07:singular-lines} shows the two parallel lines.

\begin{figure}[H]
\centering
\begin{tikzpicture}[scale=0.9]
\begin{scope}
  \draw[->,gray] (-0.4,0)--(3.4,0) node[right]{\small$x$};
  \draw[->,gray] (0,-0.4)--(0,3.4) node[above]{\small$y$};
  % line 1: x + 2y = 3 -> y = 1.5 - 0.5x
  \draw[thick,vecblue] (-0.2,1.6)--(3.2,-0.1);
  % line 2: 2x + 4y = 7 -> y = 1.75 - 0.5x
  \draw[thick,vecred] (-0.2,1.85)--(3.2,0.15);
  \node[vecred,above right] at (1.1,1.3){\small$2x{+}4y{=}7$};
  \node[vecblue,below left] at (1.7,0.55){\small$x{+}2y{=}3$};
  \node[below] at (1.5,-0.7){\small no solution (parallel)};
\end{scope}
\begin{scope}[xshift=5.8cm]
  \draw[->,gray] (-0.4,0)--(3.4,0) node[right]{\small$x$};
  \draw[->,gray] (0,-0.4)--(0,3.4) node[above]{\small$y$};
  \draw[thick,vecblue] (-0.2,1.6)--(3.2,-0.1);
  \draw[thick,vecred,dashed] (-0.2,1.6)--(3.2,-0.1);
  \draw[->,gray!70] (2.4,1.4) to[bend left=15] (1.7,0.75);
  \node[above] at (2.4,1.4){\small same line};
  \node[below] at (1.5,-0.7){\small infinitely many};
\end{scope}
\end{tikzpicture}
\caption{Two ways a $2\times 2$ system fails to have one neat solution. Left:
parallel lines that never meet, giving no solution (elimination uncovers the contradiction $0=1$). Right:
both equations represent the exact same line, giving infinitely many solutions (elimination leaves the redundant identity $0=0$, signaling a free variable).}
\label{ch07:singular-lines}
\end{figure}

\subsection{Case two: infinitely many solutions}

Now change just the right-hand side:
\[
\begin{aligned}
x + 2y &= 3\\
2x + 4y &= 6.
\end{aligned}
\qquad
\left[\begin{array}{cc|c}
1 & 2 & 3\\
2 & 4 & 6
\end{array}\right]
\;\xrightarrow{\,R_2 - 2R_1\,}\;
\left[\begin{array}{cc|c}
1 & 2 & 3\\
0 & 0 & 0
\end{array}\right].
\]
This time the last row is $0 = 0$ -- automatically true, adding no new constraint.
The two equations defined the same line (the second is just twice the first),
so every point on $x + 2y = 3$ is a solution (the right panel of
Figure~\ref{ch07:singular-lines}). There are \textbf{infinitely
many}. We say $y$ is a \emph{free variable}: it can be anything, and then $x =
3 - 2y$ follows. The solution is a whole line's worth of points.

\begin{focusbox}
A zeroed-out pivot column forces a fork in the road. Read off the right-hand
side of that dead row:
\begin{itemize}
  \item left side all zero, right side \textbf{nonzero} $\Rightarrow$ ``$0 =
  c$'' with $c\neq 0$: \textbf{no solution} (inconsistent).
  \item left side all zero, right side \textbf{zero} $\Rightarrow$ ``$0 = 0$'':
  a \textbf{free variable} and \textbf{infinitely many} solutions.
\end{itemize}
A square matrix that produces fewer than $n$ pivots is called \textbf{singular}:
it has no inverse, and the system either fails or has a whole family of answers.
A full set of $n$ nonzero pivots means $\mat{A}$ is \textbf{invertible} and the
solution is unique.
\end{focusbox}

\begin{trapbox}
Do not confuse the two failure modes. \emph{Both} produce a row whose left side
is all zeros. The distinction lies in the augmented entry on the right: a
contradictory row like ``$0=5$'' indicates \emph{no solution}; a row like
``$0=0$'' indicates a \emph{free variable}. Make sure to inspect the augmented
column before classifying the system.
\end{trapbox}

\begin{walkbox}
Take a walk and picture three planes in space. List, in your head, all the ways
they can fail to meet at one point: all three parallel; two parallel and one
crossing; three planes meeting pairwise in three parallel lines (like the three
pages near a book's spine); all three sharing a whole line; all three being the
same plane. Which of these give \emph{no} solution and which give
\emph{infinitely many}? Elimination will sort every one of these cases into the
two categories above. Notice which geometric configuration each case corresponds to.
\end{walkbox}

\section{Reduced row echelon form: the cleanest possible answer}

Upper-triangular form plus back-substitution is enough to \emph{solve} a system.
There is a tidier endpoint, though, one that makes the answer machine-readable and
gives us the natural language for the deeper structure coming in later chapters.
Instead of stopping when the matrix is triangular, we push on until it is as clean
as algebra allows. The result is the \emph{reduced row echelon
form}, almost always abbreviated \textbf{RREF}.

First, the looser notion. A matrix is in \emph{row echelon form} when:
\begin{enumerate}
  \item all-zero rows (if any) sit at the bottom, and
  \item each row's leading nonzero entry (its \emph{pivot}) is strictly to the
  right of the pivot in the row above -- so the pivots march down and to the
  right in a staircase.
\end{enumerate}
That is just upper-triangular generalized to non-square and rank-deficient
matrices. To reach the \emph{reduced} form we ask for two more things:
\begin{enumerate}
  \setcounter{enumi}{2}
  \item every pivot equals $1$ (scale each pivot row), and
  \item each pivot is the \emph{only} nonzero entry in its column -- we clear
  \emph{upward} as well as downward.
\end{enumerate}
Figure~\ref{ch07:rref-staircase} shows the resulting staircase, with leading
$1$s stepping down and to the right and free columns in between.

\begin{figure}[H]
\centering
\begin{tikzpicture}[scale=0.6]
  \draw[thick,structurepurple] (0,0) rectangle (5,4);
  \foreach \i in {1,2,3} \draw[gray!30] (0,\i)--(5,\i);
  \foreach \i in {1,2,3,4} \draw[gray!30] (\i,0)--(\i,4);
  % staircase of pivots (leading ones) marching down-right
  \fill[targetgreen!18] (0,3) rectangle (1,4);
  \fill[targetgreen!18] (1,2) rectangle (2,3);
  \fill[targetgreen!18] (3,1) rectangle (4,2);
  \node[structurepurple] at (0.5,3.5){$1$};
  \node[structurepurple] at (1.5,2.5){$1$};
  \node[structurepurple] at (3.5,1.5){$1$};
  % free / above-pivot entries marked *, zeros below
  \node at (1.5,3.5){$0$}; \node at (2.5,3.5){$*$}; \node at (3.5,3.5){$0$}; \node at (4.5,3.5){$*$};
  \node[trapred] at (0.5,2.5){$0$};
  \node at (2.5,2.5){$*$}; \node at (3.5,2.5){$0$}; \node at (4.5,2.5){$*$};
  \node[trapred] at (0.5,1.5){$0$}; \node[trapred] at (1.5,1.5){$0$};
  \node at (2.5,1.5){$0$}; \node at (4.5,1.5){$*$};
  \node[trapred] at (0.5,0.5){$0$}; \node[trapred] at (1.5,0.5){$0$};
  \node[trapred] at (2.5,0.5){$0$}; \node[trapred] at (3.5,0.5){$0$}; \node[trapred] at (4.5,0.5){$0$};
  % staircase line tracing the pivots
  \draw[thick,sourceblue,dashed] (0,4)--(1,4)--(1,3)--(2,3)--(2,2)--(4,2)--(4,1)--(5,1)--(5,0);
  \node[below] at (2.5,-0.4){\small RREF: leading $1$s (purple) step down-right; columns 3 and 5 are free};
\end{tikzpicture}
\caption{The staircase of a reduced row echelon form. Each leading $1$ (a pivot,
purple) sits strictly right of the one above; everything directly below a pivot
is forced to zero (red), and clearing upward makes each pivot the lone nonzero in
its column. Columns without a pivot ($*$ entries) belong to free variables.}
\label{ch07:rref-staircase}
\end{figure}

\begin{formalbox}
A matrix $\mat{R}$ is in \textbf{reduced row echelon form} when (i) every zero
row is below every nonzero row, (ii) the leading entry of each nonzero row is a
$1$ (a \emph{pivot} / leading $1$), (iii) each leading $1$ is to the right of the
leading $1$ above it, and (iv) each leading $1$ is the only nonzero entry in its
column. For a given matrix this form is \textbf{unique} -- you reach the same
$\mat{R}$ no matter which legal moves you choose along the way.
\end{formalbox}

Because $\mat{R}$ depends only on $\mat{A}$ and not on the elimination path, every
pivot count and free-variable pattern we read off $\mat{R}$ is an intrinsic property
of $\mat{A}$ itself.

\subsection{From triangle to RREF, traced}

Take our very first result and finish the job. We had
\[
\mat{U}=\left[\begin{array}{ccc|c}
2 & 1 & 1 & 5\\
0 & -1 & -2 & -4\\
0 & 0 & -4 & -6
\end{array}\right].
\]
\textbf{Scale each pivot to $1$.} Divide row 1 by $2$, row 2 by $-1$, row 3 by
$-4$:
\[
\left[\begin{array}{ccc|c}
1 & \tfrac12 & \tfrac12 & \tfrac52\\[2pt]
0 & 1 & 2 & 4\\[2pt]
0 & 0 & 1 & \tfrac32
\end{array}\right].
\]
\textbf{Clear upward in column 3.} Row 2 becomes (row 2) $-2\cdot$(row 3); row 1
becomes (row 1) $-\tfrac12\cdot$(row 3):
\[
\left[\begin{array}{ccc|c}
1 & \tfrac12 & 0 & \tfrac74\\[2pt]
0 & 1 & 0 & 1\\[2pt]
0 & 0 & 1 & \tfrac32
\end{array}\right].
\]
\textbf{Clear upward in column 2.} Row 1 becomes (row 1) $-\tfrac12\cdot$(row 2):
\[
\mat{R}=\left[\begin{array}{ccc|c}
1 & 0 & 0 & \tfrac54\\[2pt]
0 & 1 & 0 & 1\\[2pt]
0 & 0 & 1 & \tfrac32
\end{array}\right].
\]
In this final form, the left block is the identity matrix, and the
solution appears directly in the augmented column: $x=\tfrac54$,
$y=1$, $z=\tfrac32$ -- matching the values found by back-substitution.
\emph{RREF incorporates the back-substitution step directly into the row
operations.}

\begin{intuitionbox}
Upper-triangular form produces a system ready for back-substitution. RREF
completes those substitutions during elimination, leaving the solution directly
visible in the augmented column. The trade-off is extra arithmetic up front
(clearing upward) in exchange for an explicit readout. For solving a single
system by hand, stopping at triangular form is usually faster. For analyzing the
matrix itself -- its rank, free variables, and null space -- RREF provides the
standard canonical form used throughout subsequent chapters.
\end{intuitionbox}

\begin{labbox}
Reduce this matrix all the way to RREF by hand. It has a free variable, so the
identity will \emph{not} appear -- watch where the staircase breaks.
\[
\left[\begin{array}{ccc|c}
1 & 2 & 1 & 2\\
2 & 4 & 0 & 6\\
1 & 2 & 3 & 0
\end{array}\right].
\]
(Hint: column 2 will never get a pivot, so $y$ is free. You should find that
$z$ is pinned to a single value while $x$ depends on $y$. You will get
$z=-1$ and $x=3-2y$.)
\end{labbox}

\section{Elimination is matrix multiplication}

Here is the idea that lifts elimination from ``a procedure'' to ``a piece of
linear algebra,'' and it is the bridge to everything that follows. \textbf{Every
single row operation is the same as multiplying on the left by a special
matrix.} These special matrices are called \emph{elementary matrices}, and once
you see them, the whole algorithm becomes a product of matrices.

The trick to building an elementary matrix is almost too simple: \emph{do the row
operation to the identity matrix}, and the result is the matrix that performs
that operation on \emph{anything} when you left-multiply by it.

\subsection{The replacement matrix}

Recall our very first step: row 2 $\leftarrow$ row 2 $- 2\cdot$row 1, on a
$3$-row matrix. Apply that operation to $\mat{I}_3$. Start from
\[
\mat{I}=\begin{bmatrix}1&0&0\\0&1&0\\0&0&1\end{bmatrix},
\quad\text{do row 2}\leftarrow\text{row 2}-2\cdot\text{row 1}, \quad
\mat{E}_{21}=\begin{bmatrix}1&0&0\\-2&1&0\\0&0&1\end{bmatrix}.
\]
The only change from the identity is the $-2$ sitting in position $(2,1)$ -- the
multiplier, with a minus sign, parked exactly where the entry it kills used to
be. Now multiply $\mat{E}_{21}$ against the original augmented matrix and watch
it perform the operation:
\[
\begin{bmatrix}1&0&0\\-2&1&0\\0&0&1\end{bmatrix}
\left[\begin{array}{ccc|c}
2 & 1 & 1 & 5\\
4 & 1 & 0 & 6\\
-2 & 2 & 1 & 1
\end{array}\right]
=
\left[\begin{array}{ccc|c}
2 & 1 & 1 & 5\\
0 & -1 & -2 & -4\\
-2 & 2 & 1 & 1
\end{array}\right].
\]
Row 1 and row 3 pass through untouched (the identity rows copy them); row 2 comes
out as old-row-2 minus twice row-1. Exactly Step 1 from before. Let us verify the
new row 2 by hand: its first entry is $(-2)(2) + (1)(4) + (0)(-2) = -4 + 4 = 0$;
its second is $(-2)(1)+(1)(1) = -1$; its third $(-2)(1)+(1)(0)=-2$; the augmented
entry $(-2)(5)+(1)(6) = -4$. The $-2$ in $\mat{E}_{21}$ \emph{is} the multiplier.

\begin{focusbox}
The elementary matrix for ``row $i$ $\leftarrow$ row $i$ $-\ell\cdot$ row $j$''
is the identity with the single extra entry $-\ell$ placed at position
$(i,j)$. Left-multiplying by it performs the operation. To \emph{undo} it, flip
the sign: $\mat{E}\inv$ is the identity with $+\ell$ at $(i,j)$. (Subtracting
$\ell\cdot$row $j$ and then adding it back is the identity.)
\end{focusbox}

\begin{recoverybox}
\textbf{Pause and think:} If an elementary matrix $\mat{E}$ subtracts $3\times(\text{row } 1)$ from row $2$, what does its inverse $\mat{E}\inv$ do? You do not need any algebraic formula -- how do you undo that operation in plain words?
\end{recoverybox}

Add $3\times(\text{row } 1)$ back to row $2$. Inverting an elementary row replacement changes only the sign of the multiplier: $-3$ becomes $+3$.

\subsection{The other two as matrices}

A \textbf{scaling} of row $i$ by $c$ is the identity with $c$ in place of the
$1$ at position $(i,i)$. A \textbf{row exchange} of rows $i$ and $j$ is the
identity with rows $i$ and $j$ swapped; such a matrix is called a
\emph{permutation matrix}, written $\mat{P}$. For example, swapping rows 1 and 2
of a $3$-row matrix is done by
\[
\mat{P}=\begin{bmatrix}0&1&0\\1&0&0\\0&0&1\end{bmatrix},
\qquad
\mat{P}\begin{bmatrix}a\\b\\c\end{bmatrix}
=\begin{bmatrix}b\\a\\c\end{bmatrix}.
\]
Each of these is invertible (scaling by $c$ undoes with $1/c$; a swap undoes with
the same swap), which is the matrix version of ``every legal move is
reversible.''

\subsection{The whole elimination as one product}

Now string the steps together. Our full forward sweep was three replacement
operations: $\mat{E}_{21}$, then $\mat{E}_{31}$, then $\mat{E}_{32}$. Applying
them in order to $\mat{A}$ means left-multiplying one after another:
\[
\mat{E}_{32}\,\mat{E}_{31}\,\mat{E}_{21}\,\mat{A} = \mat{U}.
\]
Read it right to left, the order the operations actually happened. A single
matrix $\mat{E}=\mat{E}_{32}\mat{E}_{31}\mat{E}_{21}$ captures the entire forward
elimination. And because each $\mat{E}$ is invertible, we can move them to the
other side:
\[
\mat{A} = \mat{E}_{21}\inv\,\mat{E}_{31}\inv\,\mat{E}_{32}\inv\,\mat{U}
= \mat{L}\,\mat{U}.
\]
That product of inverse-elementary matrices is lower triangular, and it is the
matrix $\mat{L}$ -- whose entries are simply the multipliers
$\ell_{ij}$ themselves, sign restored. We have arrived at the
$\mat{A}=\mat{LU}$ factorization. Chapter~\ref{ch09:triangular-picture} is devoted
to exploring that factorization in full; for now, the headline is enough.

\begin{aestheticbox}
This is the moment elimination stops being arithmetic and becomes \emph{algebra}.
A sequence of row subtractions is expressed as a matrix equation
$\mat{EA}=\mat{U}$. Procedures become objects; objects can be inverted,
composed, factored, and analyzed. The viewpoint that an algebraic operation is
represented by a matrix recurs throughout the book: rotations are matrices,
projections are matrices, and derivatives on polynomial spaces are matrices.
Representing an elimination step as an elementary matrix $\mat{E}$ connects
computational procedures directly to linear algebra theory.
\end{aestheticbox}

\begin{strangbox}
``Elimination $=$ left-multiplication by elementary matrices, and undoing it
gives $\mat{A}=\mat{LU}$'' is the exact hinge on which Strang's first weeks turn.
If you are heading into his course, you are now ahead: you have already met
$\mat{E}$, $\mat{P}$, and $\mat{L}$ as concrete grids of numbers. If you are
revisiting after Strang, notice we got $\mat{L}$ for \emph{free} -- the
multipliers we were already computing \emph{are} its entries, no extra work.
That is why elimination and factorization are two views of one act.
\end{strangbox}

\subsection{Where permutations land: $\mat{PA}=\mat{LU}$}

When the elimination needs a row swap, a permutation matrix $\mat{P}$ joins the
product. Collecting all the needed swaps into a single $\mat{P}$ and applying
them first, one can show every square matrix satisfies $\mat{PA}=\mat{LU}$ for
some permutation $\mat{P}$, lower-triangular $\mat{L}$, and upper-triangular
$\mat{U}$. That is the fully general statement of what Gaussian elimination
produces, and it is exactly what a real solver computes and stores. We will prove
and exploit it in Chapter~\ref{ch09:triangular-picture}; the point here is that nothing about
row exchanges breaks the ``elimination is matrix multiplication'' story -- swaps
are just more elementary matrices.

\section{The full algorithm, in one place}

Here is the whole procedure as one recipe you can run on any system.

\begin{focusbox}
\textbf{Gaussian elimination, end to end.}
\begin{enumerate}
  \item Form the augmented matrix $[\,\mat{A}\mid\vec{b}\,]$.
  \item \emph{Forward sweep.} For each column $j$ from left to right:
  \begin{itemize}
    \item If the pivot position $(j,j)$ is zero, swap in a nonzero entry from
    below. If the whole column below is zero, move on -- this column has no
    pivot (a free variable lives here).
    \item For each row $i$ below the pivot, compute $\ell_{ij}=a_{ij}/a_{jj}$ and
    replace row $i$ by row $i$ $-\ell_{ij}\cdot$ (pivot row).
  \end{itemize}
  \item Inspect the result. A row ``$0\;\cdots\;0\mid c$'' with $c\neq 0$ means
  \emph{no solution}. Otherwise, a column without a pivot means a \emph{free
  variable} and infinitely many solutions.
  \item \emph{Solve.} Back-substitute up the triangle; or, for the cleanest
  readout, continue to RREF and read the answer off directly.
\end{enumerate}
\end{focusbox}

\subsection{One more 3x3, with a twist}

Here is a system where a swap is forced \emph{mid-stream}, not just at the start.
This is the case students most often fumble, so trace it slowly.
\[
\left[\begin{array}{ccc|c}
1 & 1 & 1 & 6\\
2 & 2 & 5 & 12\\
1 & 3 & 0 & 8
\end{array}\right].
\]
Pivot on the top-left $1$. Clear column 1. Row 2 $\leftarrow$ row 2 $-2\cdot$row
1 $= [\,0,\,0,\,3\mid 0\,]$; row 3 $\leftarrow$ row 3 $-1\cdot$row 1 $=
[\,0,\,2,\,-1\mid 2\,]$.
\[
\left[\begin{array}{ccc|c}
1 & 1 & 1 & 6\\
0 & 0 & 3 & 0\\
0 & 2 & -1 & 2
\end{array}\right].
\]
Now the second pivot position $(2,2)$ holds a $0$. But the column below it has a
$2$ in row 3 -- so we swap rows 2 and 3:
\[
\left[\begin{array}{ccc|c}
1 & 1 & 1 & 6\\
0 & 2 & -1 & 2\\
0 & 0 & 3 & 0
\end{array}\right].
\]
And we are already upper triangular -- the swap handed us a clean triangle.
Back-substitute: $3z = 0$ so $z = 0$; then $2y - 0 = 2$ so $y = 1$; then
$x + 1 + 0 = 6$ so $x = 5$. Check the original middle equation:
$2\cdot 5 + 2\cdot 1 + 5\cdot 0 = 12$. Good. Notice the pivots ended up being
$1,2,3$ -- three of them, all nonzero -- so this matrix is invertible and the
solution is the unique point $(5,1,0)$.

\begin{trapbox}
When a pivot position is zero, do \emph{not} divide by it and do \emph{not} give
up. First look down the column. Nonzero entry below? Swap. Nothing below? That
column simply has no pivot; slide your attention right to the next column while
staying in the same row, and keep going. Dividing by the zero is the single most
common fatal error -- it produces nonsense and corrupts everything downstream.
\end{trapbox}

\begin{connectionbox}
Gaussian elimination is not confined to the real numbers. The only thing the
algorithm ever does is add, subtract, multiply, and divide by nonzero pivots --
so it runs verbatim over \emph{any field}: the rationals, the complex numbers,
or the finite field of integers modulo a prime. That last case is the engine of
modern coding theory and cryptography: solving linear systems over
$\mathbb{F}_2$ (arithmetic mod $2$, where $1+1=0$) decodes error-correcting
codes and breaks linear ciphers. The schoolroom trick, run over a two-element
number system, is doing real work inside your hard drive's error correction and
your phone's signal decoder. Same algorithm, different field.
\end{connectionbox}

\section{What elimination is really telling you}

Beyond solving a single right-hand side, the number of pivots a matrix produces -- counted
across the whole forward sweep, swaps and all -- is one of the most important
numbers attached to it. It is called the \emph{rank}, and it answers questions
the original grid of numbers hid from view:

\begin{itemize}
  \item \textbf{How many genuinely independent equations are there?} Exactly the
  number of pivots. Rows that elimination flattens to zero were redundant -- they
  carried no new information.
  \item \textbf{Is the square system uniquely solvable?} Yes if and only if
  every diagonal position earns a pivot -- $n$ pivots for an $n\times n$ matrix.
  That is the same as $\mat{A}$ being invertible.
  \item \textbf{How many free variables are there?} The number of columns minus
  the number of pivots. Each free variable is a direction you can slide along
  without leaving the solution set -- the dimension of the family of answers.
\end{itemize}

Every one of these ideas -- rank, invertibility, free variables, the dimension of
the solution set -- gets its own full chapter later.

\begin{aestheticbox}
One forward sweep of row operations does far more than solve a single right-hand
side: its pivots count the independent equations, test whether $\mat{A}\inv$ exists,
count the free variables, and assemble the factorization $\mat{A}=\mat{L}\mat{U}$.
\end{aestheticbox}

Where do we go next? In this chapter, we saw that row operations correspond
to multiplications on the left by elementary matrices $\mat{E}$, and inverting those
operations begins to assemble the matrix factorization $\mat{A}=\mat{L}\mat{U}$.
In the next chapter, we examine matrix multiplication and inverses more broadly.
What does it mean to multiply one matrix by another as a unified geometric transformation?
Why does the row-by-column rule work the way it does? And when a square matrix has a
full set of $n$ pivots, how do we find a single inverse matrix $\mat{A}\inv$ that
undoes $\mat{A}$ directly, turning $\mat{A}\vec{x}=\vec{b}$ into
$\vec{x}=\mat{A}\inv\vec{b}$? The next chapter, \emph{Matrix Multiplication and Inverses},
builds that algebraic foundation.

\begin{historybox}
Although named after Carl Friedrich Gauss, systematic elimination began much earlier.
Around 150~BCE, Chinese scholars compiling \emph{The Nine Chapters on the Mathematical Art}
solved agricultural tax and grain-yield problems involving several simultaneous conditions.
Working with bamboo counting rods arranged in rectangular arrays on a grid board (\emph{Fangcheng}),
they repeatedly subtracted multiples of one column from another to clear entries below the diagonal --
performing the forward sweep of elimination over eighteen centuries before Europe codified it.

In 1801, Gaussian elimination took its modern shape through celestial mechanics. On New Year's Day,
the Italian astronomer Giuseppe Piazzi spotted a dim, wandering body in the night sky -- the dwarf
planet Ceres. Piazzi tracked it for only 41 days across a tiny nine-degree arc before it was lost
in the blinding glare of the Sun. Astronomers across Europe despaired of finding it again: they had
only a handful of noisy, imperfect observations, and Newton's gravitational laws led to non-linear
orbital equations that were impossible to solve directly.

The 24-year-old Carl Friedrich Gauss took up the challenge. Developing the \emph{method of least
squares} -- a technique for finding the single best-fitting mathematical curve through a cloud of noisy,
imperfect observations -- Gauss set up a system of six coupled linear equations in six unknown orbital parameters.
With only pen, paper, and logarithms as tools, he organized the arithmetic into systematic forward elimination
and backward substitution, cutting the hand computation down to manageable hours. On New Year's Eve 1801,
guided by Gauss's calculations, astronomers pointed their telescopes into the night sky and found Ceres
within half a degree of his predicted position. Gauss later published the systematic algorithm in 1809
in \emph{Theoria Motus Corporum Coelestium} while calculating the orbit of the asteroid Pallas.

Eighty years later, in 1888, the German geodesist Wilhelm Jordan was tackling \emph{geodesy} -- the science
of accurately measuring the Earth's geometric shape and mapping regional networks of surveying triangles. Facing
hundreds of error-prone equations, Jordan extended Gauss's procedure by using each pivot to eliminate entries
\emph{upward} as well as downward. This drove the matrix straight to reduced row echelon form (RREF) or the
identity matrix, allowing surveyors to read off coordinates directly without a separate back-substitution
climb. The combined algorithm has been known as \emph{Gauss--Jordan elimination} ever since.

In pure mathematics with exact fractions, elimination succeeds whenever pivots are nonzero. But on
digital computers working in \emph{floating-point arithmetic} -- where numbers are stored with a finite
number of digits and rounded off after every operation -- naive elimination is dangerous. If a pivot is
tiny (say $10^{-5}$), dividing by it creates enormous multipliers ($10^5$), and subsequent row subtractions
can magnify microscopic roundoff errors until they completely swamp the true solution.

In 1947, computing pioneers John von Neumann and Herman Goldstine published a paper warning that
Gaussian elimination on digital computers would suffer catastrophic, exponential roundoff growth,
predicting that errors would explode on matrices with more than 15 or 20 variables. That sparked an urgent
controversy over whether linear algebra could ever run reliably on electronic computers. Alan Turing
(in his 1948 paper \emph{Rounding-Off Errors in Matrix Processes}) and later James H. Wilkinson (in 1961)
resolved the crisis using \emph{backward error analysis} -- a clever shift in perspective that evaluates
stability by asking not ``how wrong is my answer for the original problem?'' but ``what slightly perturbed
problem did my computed answer solve exactly?'' They proved that with \textbf{partial pivoting} -- swapping
rows at each step to seat the largest available entry as the pivot -- multipliers are kept within
$|\ell_{ij}| \le 1$, making Gaussian elimination numerically stable for almost all practical matrices.

Gaussian elimination provided science with its first universal, $O(n^3)$ algorithm: it solves systems
of any size, checks whether solutions exist, tracks families of solutions when free variables appear,
and measures the matrix's rank. Yet two major frontiers remain open today:
\begin{itemize}
  \item \textbf{The optimal complexity exponent:} Is solving an $n\times n$ system fundamentally an
  $O(n^3)$ problem, or can it be done essentially in $O(n^2)$ time -- as fast as reading the $n^2$
  entries of the matrix? In 1969, Volker Strassen showed that matrix multiplication
  (and thus elimination) can be carried out in $O(n^{2.81})$ operations. Theoretical computer scientists
  have pushed that theoretical exponent down near $O(n^{2.37155})$, but the true optimal exponent $\omega$ remains
  one of the most famous open problems in computer science.
  \item \textbf{Worst-case stability vs.\ speed:} Although partial pivoting runs in cheap $O(n^2)$ extra
  time and never fails on practical problems, artificial matrices exist (such as Kahan's matrix) where its roundoff
  theoretically blows up exponentially. \emph{Complete pivoting} (swapping both rows and columns) guarantees
  provable stability for every matrix, but searching the entire submatrix costs an expensive $O(n^3)$ overhead.
  Finding a pivoting strategy that combines the speed of partial pivoting with stronger theoretical
  stability guarantees remains an active research direction in numerical linear algebra.
\end{itemize}

There is another practical bottleneck: if you solve $\mat{A}\vec{x}=\vec{b}$ and subsequently change only the right-hand side $\vec{b}$ (such
as testing multiple wind loads on the same bridge structure), running naive elimination from scratch wastes
$O(n^3)$ work by repeating the exact same row operations on $\mat{A}$. Overcoming this repetitive cost
requires packaging elimination into matrix operators: Chapter~\ref{ch08:origin} examines matrix
multiplication and inverses, paving the way for Chapter~\ref{ch09:triangular-picture} to introduce
LU factorization ($\mat{A} = \mat{L}\mat{U}$), which stores the multipliers so that
any new right-hand side can be solved in $O(n^2)$ time.
\end{historybox}

\section{Exercises}

\begin{exercise}[Mental checks: true or false?]
Decide whether each statement is true or false, and give a one-sentence geometric or algebraic reason:
\begin{enumerate}[label=(\alph*)]
  \item Multiplying an entire equation by a nonzero constant changes the line or plane it represents in space.
  \item In an upper-triangular system $\mat{U}\vec{x} = \vec{c}$, back-substitution solves the system without computing any new multipliers $\ell_{ij}$.
  \item If an $n \times n$ system produces fewer than $n$ nonzero pivots during elimination, it is guaranteed to have no solution.
  \item Swapping two rows of an augmented matrix changes the coordinates $(x, y, z)$ of the intersection point.
\end{enumerate}
\end{exercise}

\begin{exercise}[Warm-up, fully traced]
Solve by Gaussian elimination, showing the multiplier at every step and then
back-substituting:
\[
\begin{aligned}
x + 2y + z &= 4\\
2x + 5y + 3z &= 11\\
2x + 3y + 2z &= 8.
\end{aligned}
\]
Verify your answer in all three original equations.
\end{exercise}

\begin{exercise}[A forced swap]
Run elimination on
\[
\left[\begin{array}{ccc|c}
0 & 1 & 2 & 3\\
1 & 2 & 1 & 4\\
3 & 1 & 0 & 5
\end{array}\right].
\]
You must exchange rows before you can begin. Identify which swap, reach
upper-triangular form, and back-substitute. State the three pivots.
\end{exercise}

\begin{exercise}[Diagnose the failure]
For each system, run elimination far enough to decide whether it has no solution,
exactly one, or infinitely many. If infinitely many, describe the solution set
(name the free variable and write the others in terms of it).
\[
\text{(a)}\;\;
\begin{aligned}
x + y &= 2\\
3x + 3y &= 5,
\end{aligned}
\qquad
\text{(b)}\;\;
\begin{aligned}
x - y &= 1\\
-2x + 2y &= -2.
\end{aligned}
\]
\end{exercise}

\begin{exercise}[All the way to RREF]
Take the upper-triangular matrix
\[
\left[\begin{array}{ccc|c}
3 & 6 & 1 & 0\\
0 & 2 & 4 & 8\\
0 & 0 & 5 & 10
\end{array}\right]
\]
and reduce it to RREF: scale each pivot to $1$ and clear upward. Read off
$x,y,z$ from the final augmented column.
\end{exercise}

\begin{exercise}[Elementary matrices by hand]
Write down the $3\times 3$ elementary matrix $\mat{E}$ that performs ``row 3
$\leftarrow$ row 3 $-4\cdot$ row 1.'' Multiply it against
$\mat{A}=\colvec{1&0&2\\0&1&1\\4&2&0}$ and confirm by direct multiplication that
the product has a zero in position $(3,1)$. Then write $\mat{E}\inv$ and verify
$\mat{E}\inv\mat{E}=\mat{I}$.
\end{exercise}

\begin{exercise}[A free variable in 3D]
Solve
\[
\begin{aligned}
x + y + z &= 3\\
2x + 2y + 3z &= 8\\
x + y + 2z &= 5,
\end{aligned}
\]
and notice that column 2 never produces a pivot. Identify the free variable and
write the full family of solutions as a point plus a multiple of a direction
vector. (Geometrically, the three planes share a whole line.)
\end{exercise}

\begin{exercise}[Tiny-pivot instability]
Consider, with exact arithmetic,
\[
\begin{aligned}
0.001\,x + y &= 1\\
x + y &= 2.
\end{aligned}
\]
(a) Solve exactly. (b) Now pretend your calculator keeps only $3$ significant
digits and eliminate \emph{without} swapping -- using $0.001$ as the pivot.
Watch the answer for $x$ degrade. (c) Swap the two rows first (partial pivoting)
and redo with $3$ digits. Explain why the swapped version is far more accurate.
\end{exercise}

\begin{exercise}[Counting operations]
Argue that clearing column $1$ of an $n\times n$ matrix costs about $n^2$
multiply-add operations (you touch each of the $n-1$ rows below the pivot, and
each touched row has about $n$ entries). Summing the shrinking columns, explain
informally why the whole forward sweep costs on the order of $\tfrac13 n^3$.
(Hint: $\sum_{k=1}^{n} k^2 \approx \tfrac13 n^3$.)
\end{exercise}

\begin{exercise}[Uniqueness of RREF -- a proof]
Suppose two different sequences of legal row operations turn the same matrix
$\mat{A}$ into row echelon forms $\mat{R}_1$ and $\mat{R}_2$. Argue that
$\mat{R}_1$ and $\mat{R}_2$ have their pivots in exactly the same columns.
(Hint: a column is a pivot column precisely when it is \emph{not} a linear
combination of the columns to its left, and row operations do not change which
columns are combinations of which earlier columns. You may take this last fact
as given for now -- we prove it carefully in the chapter on the column space.)
\end{exercise}

\begin{exercise}[Elimination over a tiny field]
Work over $\mathbb{F}_2$, where the only numbers are $0$ and $1$ and $1+1=0$.
Solve
\[
\begin{aligned}
x + y \phantom{{}+ z} &= 1\\
x \phantom{{}+ y} + z &= 0\\
y + z &= 1,
\end{aligned}
\]
using exactly the same elimination, but with mod-$2$ arithmetic (so subtracting
equals adding). Find all solutions. This provides a simple illustration of decoding an
error-correcting code.
\end{exercise}
